\chapter{课堂录音进入公开写作的处理流程}

课堂原声能让教学写作摆脱空泛，但也带来转写误差、隐私与选择性引用风险。建议采用以下流程。

\begin{enumerate}
  \item \textbf{归档：}保留原始音频、转写文件、日期、时长和来源平台；
  \item \textbf{定位：}每条候选引文记录文件名和时间戳；
  \item \textbf{回听：}正式出版前回听关键句，修正自动转写错误；
  \item \textbf{去识别：}删除姓名、班级、账号、个人经历等非必要信息；
  \item \textbf{取授权：}公开学生作品、长篇发言或可识别经历前取得授权；
  \item \textbf{保语境：}同时记录前后问题，避免把犹豫、玩笑和现场估计写成定论；
  \item \textbf{分证据：}教师分享属于实践反思，学生发言属于个案材料，均不自动成为效果数据；
  \item \textbf{建台账：}对每条公开引文记录状态：待回听、已回听、已授权、仅内部使用。
\end{enumerate}

\begin{tech-note}
公开 GitHub 仓库不应上传原始课堂录音和未经匿名化的完整转写。书中可保留日期与时间戳，原始证据放在受控的私有档案中。
\end{tech-note}
